\chapter{Tensors and Exactness}

\begin{center}\itshape
The quotient map $\Z\to\Z/2$ remembers parity and forgets even numbers. An exact sequence records both facts. What happens to that record when we apply $\operatorname{Hom}$ or take a tensor product?
\end{center}

Start with the sequence
\[
0\longrightarrow\Z\xrightarrow{\,2\,}\Z
\longrightarrow\Z/2\longrightarrow0.
\]
The first map sends $a$ to $2a$; the second takes the residue class modulo $2$. Every residue class has a representative. The representatives of $0$ are exactly the even integers, and each even integer has a unique preimage under the first map. The sequence puts these three assertions on one line.

Now try to choose representatives by a homomorphism $s\colon\Z/2\to\Z$. If $s$ were a section of the quotient map, $s(\bar1)$ would be odd. But $2s(\bar1)=s(\bar0)=0$, and the only integer killed by $2$ is $0$. Such a choice is impossible. Exactness records the kernel and the quotient; it does not say that the middle module is their direct sum.

Chapter~\ref{linear:ch:tensors} built tensor products of vector spaces from bases. Modules bring a new difficulty: they need not have bases, and equations such as $2x=0$ can hold for non-zero elements. We therefore define the tensor product through its maps, using the universal properties of Chapter~1. Its effect on exact sequences will tell us which part of the record it preserves.

There is a memorable picture of the other main construction of this chapter. In the opening of \emph{It's My Turn} (1980), Jill Clayburgh plays a mathematics professor proving the snake lemma at a blackboard. Benedict Gross advised on the mathematics and coached her in the lines. The lemma builds a map from two exact rows of a diagram. We will follow an element through that diagram and check why the answer does not depend on the choices made along the way. The same construction underlies the connecting maps in the Mayer--Vietoris sequence of Volume~9; Chapter~7 will use it for resolutions of modules.

Hassler Whitney treated tensor products of abelian groups in 1938. That case is a useful guide throughout: an abelian group is a module over $\Z$. Unless another ring is specified, $R$ is a commutative ring with identity, and all modules and maps are unital $R$-modules and $R$-linear maps. The module language was introduced in Volume~3, Chapter~7.

\section{Exact sequences}

An element that disappears under one map should come from the preceding map. This is the condition called exactness.

\begin{definition}\label{def:exact}
For consecutive homomorphisms
\[
\cdots\longrightarrow M_{i-1}\xrightarrow{f_{i-1}}M_i
\xrightarrow{f_i}M_{i+1}\longrightarrow\cdots,
\]
\emph{exactness at $M_i$} means $\operatorname{im}f_{i-1}=\ker f_i$. The sequence is \emph{exact} if this holds at each term that has both a preceding and a following map. A \emph{short exact sequence} has the form
\[
0\longrightarrow A\xrightarrow f B\xrightarrow g C\longrightarrow0.
\]
Thus $f$ is injective, $g$ is surjective, and $f(A)=\ker g$.
\end{definition}

The zero terms have a job. Exactness at $A$ says $\ker f=0$, because the image of the map from $0$ is $0$. Exactness at $C$ says $\operatorname{im}g=C$, because the kernel of the map to $0$ is all of $C$. The first isomorphism theorem then identifies $C$ with $B/f(A)$.

For a map $h\colon M\to N$, its \emph{cokernel} is $N/\operatorname{im}h$. It records the classes in $N$ that $h$ does not reach. The definition gives the exact sequence
\[
0\to\ker h\to M\xrightarrow h N\to\operatorname{coker}h\to0.
\]
The middle two arrows alone already describe every failure of $h$ to be an isomorphism: an isomorphism has both kernel and cokernel zero, and the converse follows from injectivity and surjectivity.

There is always a short exact sequence with middle term $A\oplus C$: use $a\mapsto(a,0)$ and $(a,c)\mapsto c$. A short exact sequence is called \emph{split} when it is isomorphic to this one, with the maps from $A$ and to $C$ respected.

\begin{example}\label{ex:ses}
Neither
\[
0\to\Z\xrightarrow{\,2\,}\Z\to\Z/2\to0
\quad\text{nor}\quad
0\to\Z/2\xrightarrow{\,2\,}\Z/4\to\Z/2\to0
\]
splits. In the second sequence the first map sends $\bar a$ modulo $2$ to $\overline{2a}$ modulo $4$. A split first sequence would give $\Z\cong\Z\oplus\Z/2$, although the right side has a non-zero element of finite order and the left side does not. A split second sequence would give $\Z/4\cong\Z/2\oplus\Z/2$, although only the left side has an element of order $4$. In particular, the two end modules do not determine the middle one.
\end{example}

A splitting can be recognised by choosing a section of the quotient map, or by retracting the middle module onto the kernel. The next result makes both tests precise.

\begin{proposition}[splitting lemma]\label{prop:split}
Given $0\to A\xrightarrow f B\xrightarrow g C\to0$, the following conditions are equivalent.
\begin{enumerate}[label=(\alph*)]
\item A homomorphism $r\colon B\to A$ satisfies $rf=\mathrm{id}_A$. Such a map is a \emph{retraction} of $f$.
\item A homomorphism $s\colon C\to B$ satisfies $gs=\mathrm{id}_C$. Such a map is a \emph{section} of $g$.
\item An isomorphism $\varphi\colon B\to A\oplus C$ satisfies $\varphi(f(a))=(a,0)$ and $p_C\varphi=g$, where $p_C(a,c)=c$.
\end{enumerate}
\end{proposition}

\begin{proof}
\emph{A retraction gives coordinates.} Under (a), put $\varphi(b)=(r(b),g(b))$. The equation $gf=0$ gives $\varphi(f(a))=(a,0)$. If both coordinates of $\varphi(b)$ vanish, exactness gives $b=f(a)$, and then $a=r(f(a))=r(b)=0$. Thus $\varphi$ is injective.

To obtain a prescribed pair $(a,c)$, first choose $b_0$ with $g(b_0)=c$. It may have the wrong first coordinate. Correct it inside the kernel by setting
\[
b=b_0+f(a-r(b_0)).
\]
Then $g(b)=c$ and $r(b)=a$. This proves surjectivity and (c).

\emph{A section separates the two pieces.} Under (b), define $\psi(a,c)=f(a)+s(c)$. If this is $0$, apply $g$ to get $c=0$, and then use injectivity of $f$ to get $a=0$. For any $b\in B$, the difference $b-s(g(b))$ belongs to $\ker g=f(A)$. It follows that $b=\psi(a,g(b))$ for some $a$. So $\psi$ is an isomorphism, and its inverse has the properties in (c).

\emph{Coordinates give both maps back.} Under (c), take $r=p_A\varphi$ and $s(c)=\varphi^{-1}(0,c)$. These satisfy $rf=\mathrm{id}_A$ and $gs=\mathrm{id}_C$.
\end{proof}

For a vector space quotient, a basis of the quotient can be lifted and extended linearly, giving a section. The example with $\Z/2$ shows exactly where this argument can fail for modules: the lifted elements must also obey the relations of the quotient.

\section{The snake lemma}

Suppose that a map between two exact sequences is given. Its three component maps have kernels and cokernels of their own. The snake lemma joins these six modules, with one new map between them.

\begin{theorem}[snake lemma]\label{thm:snake}
Let the rows $A\xrightarrow f B\xrightarrow g C\to0$ and $0\to A'\xrightarrow{f'}B'\xrightarrow{g'}C'$ be exact. Let $a\colon A\to A'$, $b\colon B\to B'$, and $c\colon C\to C'$ satisfy $bf=f'a$ and $cg=g'b$, as in Figure~\ref{fig:snake}. There is an $R$-linear \emph{connecting homomorphism} $\partial\colon\ker c\to\operatorname{coker}a$ for which
\[
\begin{gathered}
\ker a\xrightarrow f\ker b\xrightarrow g\ker c
\xrightarrow{\partial}\operatorname{coker}a\\
\xrightarrow{\bar f'}\operatorname{coker}b
\xrightarrow{\bar g'}\operatorname{coker}c
\end{gathered}
\]
is exact. Here $f,g$ are restricted to kernels, and $\bar f',\bar g'$ are induced on cokernels. If $f$ is injective, the first arrow is injective. If $g'$ is surjective, the last arrow is surjective.
\end{theorem}

The two exactness assumptions have different ends: $g$ must be onto, and $f'$ must be injective. We need the first to lift an element of $C$ to $B$, and the second to pull an element of $f'(A')$ back to a unique element of $A'$.

\begin{figure}[ht]
\centering
\fig{ch02-snake}
\caption{The snake lemma. The kernels of the vertical maps sit above the diagram, their cokernels below; the connecting homomorphism $\partial$ is the snake's body, running from the
top right to the bottom left. To compute $\partial(z)$: lift $z$ to $y\in B$, push it down to $b(y)\in B'$, which lies in $f'(A')$, pull it back to $x'\in A'$, and take its class.}
\label{fig:snake}
\end{figure}

\begin{proof}
\emph{Check the four inherited arrows.} For $x\in\ker a$, the equation $b(f(x))=f'(a(x))=0$ puts $f(x)$ in $\ker b$. The equation $c(g(y))=g'(b(y))$ does the same for $g$. Also $f'(a(x))=b(f(x))$ lies in $\operatorname{im}b$. Hence $f'$ sends $\operatorname{im}a$ into $\operatorname{im}b$ and descends to the quotient map $\bar f'$. The second square gives $\bar g'$ in the same way.

\emph{Construct the missing arrow.} Take $z\in\ker c$. Choose $y\in B$ with $g(y)=z$. Commutativity gives
\[
g'(b(y))=c(g(y))=c(z)=0.
\]
Exactness of the bottom row puts $b(y)$ in $\operatorname{im}f'$, so there is a unique $x'\in A'$ with $f'(x')=b(y)$. Define
\[
\partial(z)=x'+\operatorname{im}a.
\]
This quotient removes the ambiguity in choosing $y$. Indeed, another lift is $y+f(x)$ for some $x\in A$, because $\ker g=\operatorname{im}f$. Its image under $b$ is $f'(x'+a(x))$, so it changes $x'$ only by an element of $\operatorname{im}a$. The class is unchanged. To check addition, use $y_1+y_2$ as a lift of $z_1+z_2$; its pullback is $x'_1+x'_2$. To check multiplication by $r\in R$, use $ry_1$ and $rx'_1$. Thus $\partial$ is linear.

\emph{Exactness at $\ker b$.} A value $f(x)$ with $a(x)=0$ lies in $\ker b$ and is killed by $g$. Conversely, let $y\in\ker b$ with $g(y)=0$. Write $y=f(x)$. Then $f'(a(x))=b(y)=0$, and injectivity of $f'$ gives $a(x)=0$.

\emph{Exactness at $\ker c$.} If $z=g(y)$ for $y\in\ker b$, use that lift in the construction. Its image is $b(y)=0=f'(0)$, so $\partial(z)=0$. Conversely, if $\partial(z)=0$, the constructed $x'$ equals $a(x)$ for some $x$. Replacing $y$ by $y-f(x)$ leaves $g(y)$ unchanged and makes its image under $b$ equal to $f'(x'-a(x))=0$. Thus $z$ comes from $\ker b$.

\emph{Exactness at $\operatorname{coker}a$.} The class $\partial(z)=[x']$ maps to $[f'(x')]=[b(y)]=0$ in $\operatorname{coker}b$. Conversely, if a class $[x']$ maps to zero there, write $f'(x')=b(y)$. Put $z=g(y)$. Then
\[
c(z)=g'(b(y))=g'(f'(x'))=0,
\]
and this choice of lift gives $\partial(z)=[x']$.

\emph{Exactness at $\operatorname{coker}b$.} The composite $\bar g'\bar f'$ vanishes because $g'f'=0$. If $[y']$ maps to zero under $\bar g'$, then $g'(y')=c(z)$ for some $z\in C$. Lift $z=g(y)$. Now $g'(y'-b(y))=0$, so $y'-b(y)=f'(x')$ for some $x'$. Modulo $\operatorname{im}b$, this says $[y']=\bar f'([x'])$.

Finally, a restriction of an injective $f$ is injective. If $g'$ is onto, every class $[z']$ has a representative $z'=g'(y')$, so $\bar g'$ is onto as well. No condition on $f$ or $g'$ was needed for the four inner exactness checks.
\end{proof}

The zero element causes no exception in the construction: lift it by $0$, then pull back $0$, and obtain the zero class. The proof also explains the direction of $\partial$. The construction starts in $\ker c$, at the upper right, and ends in $\operatorname{coker}a$, at the lower left.

\begin{corollary}[short five lemma]\label{cor:fivelemma}
Suppose $a,b,c$ form a commutative diagram between short exact sequences $0\to A\to B\to C\to0$ and $0\to A'\to B'\to C'\to0$. If $a$ and $c$ are isomorphisms, then $b$ is an isomorphism.
\end{corollary}

\begin{proof}
The snake lemma gives exact sequences
\[
0\to\ker a\to\ker b\to\ker c,
\qquad
\operatorname{coker}a\to\operatorname{coker}b
\to\operatorname{coker}c\to0.
\]
Both outside kernels are zero, which forces $\ker b=0$. Both outside cokernels are zero, which forces $\operatorname{coker}b=0$. Thus $b$ is injective and surjective.
\end{proof}

\begin{think}
Take two short exact sequences with the same ends $A,C$. A map between their middle terms that respects the inclusions of $A$ and the projections to $C$ has identity maps at both ends. Apply the corollary: why is the middle map automatically invertible? Chapter~7 uses this fact when comparing extensions.
\end{think}


% V10-EXPANSION WORKED EXAMPLE 1
\subsection*{Worked example: a nonzero connecting map}

Use the same short exact sequence in both rows of a snake diagram:
\[
0\longrightarrow\Z\xrightarrow{\,2\,}\Z
 \longrightarrow\Z/2\longrightarrow0.
\]
Let the first two vertical maps be multiplication by $2$, and let the last vertical map on $\Z/2$ be zero. The left square commutes because $2\cdot2=2\cdot2$. The right square commutes because reducing $2k$ modulo $2$ gives zero. The construction is therefore within the hypotheses of the snake lemma.

The first two vertical maps have zero kernels and cokernel $\Z/2$. The last has both kernel and cokernel $\Z/2$. Thus the new arrow is
\[
\partial:\Z/2\longrightarrow\Z/2.
\]
To find it, start with $\bar1$ in the last kernel and lift it to $1\in\Z$ in the upper middle term. The middle vertical map sends this to $2$. Pulling $2$ back through the lower injection, also multiplication by $2$, gives $1$ in the lower left term. Its class in the first cokernel is $\bar1$. Hence $\partial$ is the identity.

Check the choice in this calculation. Any other lift is $1+2t$. Its downward image is $2+4t$, whose pullback is $1+2t$; its class modulo $2$ is still $\bar1$. The quotient in the target removes precisely the dependence on the lift.

After deleting the first zero kernels, the resulting sequence reads
\[
0\longrightarrow\Z/2\xrightarrow{\,\mathrm{id}\,}\Z/2
 \xrightarrow{\,0\,}\Z/2
 \xrightarrow{\,\mathrm{id}\,}\Z/2\longrightarrow0.
\]
The middle zero arrow is induced by the lower injection $\Z\xrightarrow2\Z$ on the two cokernels. It is zero because multiplication by $2$ is zero modulo $2$. The final identity is induced by reduction modulo $2$. Exactness is now visible at every term. A diagram whose vertical maps look very simple can have a nonzero connecting map; the map measures how lifting and moving between the rows interact.

\section{Hom and exactness}

Fix $M$. Applying $\operatorname{Hom}(M,-)$ replaces a module by all the maps into it from $M$. We can add those maps pointwise and multiply them by elements of $R$:
\[
(u+v)(x)=u(x)+v(x),\qquad (ru)(x)=r\,u(x).
\]
Commutativity of $R$ ensures that $ru$ is still $R$-linear. A map $h\colon N\to N'$ acts by postcomposition $h_*(u)=h\circ u$. If instead we fix the target $N$, a map $h\colon M'\to M$ acts by precomposition $h^*(u)=u\circ h$. These are the covariant and contravariant functors of Example~\ref{ex:functors}(f).

\begin{proposition}[Hom is left exact]\label{prop:homleft}
\leavevmode
\begin{enumerate}[label=(\alph*)]
\item Exactness of $0\to A\xrightarrow f B\xrightarrow g C$ implies exactness of
\[
0\to\operatorname{Hom}(M,A)\xrightarrow{f_*}\operatorname{Hom}(M,B)
\xrightarrow{g_*}\operatorname{Hom}(M,C).
\]
\item Exactness of $A\xrightarrow f B\xrightarrow g C\to0$ implies exactness of
\[
0\to\operatorname{Hom}(C,N)\xrightarrow{g^*}\operatorname{Hom}(B,N)
\xrightarrow{f^*}\operatorname{Hom}(A,N).
\]
\end{enumerate}
\end{proposition}

\begin{proof}
\emph{Maps into a kernel.} In (a), $fu=0$ implies $u=0$ because $f$ is injective. Also $gfu=0$ for every $u$. Suppose $v\colon M\to B$ satisfies $gv=0$. Each $v(m)$ lies in $f(A)$, so there is a unique $u(m)\in A$ with $f(u(m))=v(m)$. The uniqueness checks linearity: applying $f$ to $u(m+m')$ and to $u(m)+u(m')$ gives the same value, and similarly for scalar multiples. Hence $v=fu$ with $u$ linear.

\emph{Maps out of a quotient.} In (b), $ug=0$ implies $u=0$ because $g$ is onto, and every $ug$ vanishes on $f(A)$. Conversely, if $v\colon B\to N$ vanishes on $f(A)=\ker g$, define $u(c)=v(b)$ using any lift $g(b)=c$. Two lifts differ by an element of $\ker g$, so the value is independent of the lift. Addition and scalar multiplication can be checked on lifts. Thus $u$ is linear and $v=ug$.
\end{proof}

Neither assertion promises surjectivity at the right end. For example,
$\operatorname{Hom}(\Z/2,\Z)=0$: the image of $\bar1$ would have to be killed by $2$. Therefore the identity map of $\Z/2$ cannot lift along $\Z\to\Z/2$. Applying $\operatorname{Hom}(\Z/2,-)$ to the opening short exact sequence loses its last surjection.

We will also need a way to read exactness back from the maps into every possible target.

\begin{lemma}\label{lem:testexact}
Let $A\xrightarrow f B\xrightarrow g C$ be module maps. If, for every module $N$,
\[
0\to\operatorname{Hom}(C,N)\xrightarrow{g^*}\operatorname{Hom}(B,N)
\xrightarrow{f^*}\operatorname{Hom}(A,N)
\]
is exact, then $A\xrightarrow f B\xrightarrow g C\to0$ is exact.
\end{lemma}

\begin{proof}
We choose three targets, one for each assertion.

For surjectivity of $g$, use $N=C/\operatorname{im}g$. The quotient map $\pi\colon C\to N$ satisfies $\pi g=0$. Injectivity of $g^*$ gives $\pi=0$, so this quotient is zero and $g$ is onto.

For $gf=0$, use $N=C$. Exactness makes $f^*g^*=0$; applying this composite to $\mathrm{id}_C$ gives $gf=0$. Thus $\operatorname{im}f\subseteq\ker g$.

For the reverse inclusion, use $N=B/\operatorname{im}f$. Its quotient map $\rho$ satisfies $\rho f=0$. Exactness provides $u\colon C\to N$ with $\rho=ug$. Consequently $g(b)=0$ implies $\rho(b)=0$, or $b\in\operatorname{im}f$. The two inclusions prove exactness at $B$.
\end{proof}

\section{The tensor product}

For integers, the map $(a,b)\mapsto ab$ is linear in each argument separately, although it is not linear as a map from the direct sum $\Z\oplus\Z$ to $\Z$. For instance, $(1,0)$ and $(0,1)$ both have product $0$, whereas their sum $(1,1)$ has product $1$. A map $\beta\colon M\times N\to P$ with separate linearity is called \emph{bilinear}. Other examples are a dot product, evaluation $V^*\times V\to K$, and the determinant of a $2\times2$ matrix viewed as a function of its two columns. A bilinear map need not be linear on $M\oplus N$.

We want a module in which the symbol $m\otimes n$ obeys precisely the relations that every bilinear map obeys. A linear map from that module should then be the same thing as a bilinear map from $M\times N$.

\begin{definition}\label{def:tensor}
A \emph{tensor product} of $M,N$ over $R$ consists of a module $M\otimes_RN$ and a bilinear map $(m,n)\mapsto m\otimes n$ with the following property: for every module $P$ and every bilinear $\beta\colon M\times N\to P$, there is a unique linear map $\tilde\beta\colon M\otimes_RN\to P$ satisfying
\[
\tilde\beta(m\otimes n)=\beta(m,n).
\]
\end{definition}

The uniqueness matters as much as existence: agreement on all symbols $m\otimes n$ must determine the entire linear map.

\begin{theorem}\label{thm:tensorexists}
For any modules $M,N$, a tensor product exists. Any two are related by a unique isomorphism respecting their bilinear maps. The elements $m\otimes n$, called \emph{pure tensors}, generate $M\otimes_RN$.
\end{theorem}

\begin{proof}
\emph{Build a module with the required relations.} Let $F$ be free on the set $M\times N$. Its basis vector indexed by $(m,n)$ will be written $[m,n]$. Let $K$ be generated by
\[
\begin{gathered}
[m+m',n]-[m,n]-[m',n],\\
[m,n+n']-[m,n]-[m,n'],\\
[rm,n]-r[m,n],\qquad [m,rn]-r[m,n].
\end{gathered}
\]
In $F/K$, define $m\otimes n=[m,n]+K$. Each displayed relation becomes zero, so this map is bilinear.

Given $\beta$, define the linear map $F\to P$ by $[m,n]\mapsto\beta(m,n)$. Bilinearity says that it kills each generator of $K$, hence all of $K$. It descends to the desired $\tilde\beta\colon F/K\to P$. Its values on the pure tensors are prescribed, and those tensors generate the quotient. This proves uniqueness of $\tilde\beta$ and the generation assertion.

\emph{Compare two constructions.} If $T,T'$ both have the defining property, it gives a map $u\colon T\to T'$ sending each $m\otimes n$ to the corresponding tensor in $T'$, and a map $v$ in the other direction. The composites agree with the identities on all pure tensors. Uniqueness makes $vu=\mathrm{id}_T$ and $uv=\mathrm{id}_{T'}$. It also makes $u$ the only isomorphism with this compatibility. This is the uniqueness of initial objects from Proposition~\ref{prop:initialunique}, written out for bilinear maps.
\end{proof}

An element of the tensor product is a finite sum of pure tensors; it need not itself be pure. More surprisingly, a pure tensor of two non-zero elements can be zero. In $\Z/2\otimes_\Z\Z/3$, for example,
\[
x\otimes y=(3x-2x)\otimes y
=x\otimes3y-(2x)\otimes y=0.
\]
The relation $rm\otimes n=m\otimes rn$ transfers a scalar from one side to the other, where it may kill the element.

If $u\colon M\to M'$ and $v\colon N\to N'$ are linear, the map $(m,n)\mapsto u(m)\otimes v(n)$ is bilinear. The defining property therefore gives
\[
u\otimes v\colon M\otimes N\to M'\otimes N',\qquad
m\otimes n\mapsto u(m)\otimes v(n).
\]
Checking on pure tensors shows that identities and composites are respected. In particular, tensoring with a fixed $N$ is a functor.

\begin{proposition}\label{prop:tensorrules}
The following maps are natural isomorphisms.
\begin{enumerate}[label=(\alph*)]
\item $R\otimes_R M\to M$, given by $r\otimes m\mapsto rm$.
\item $M\otimes_RN\to N\otimes_RM$, given by $m\otimes n\mapsto n\otimes m$.
\item $(\bigoplus_iM_i)\otimes_RN\to\bigoplus_i(M_i\otimes_RN)$, given by $(m_i)_i\otimes n\mapsto(m_i\otimes n)_i$.
\item For an ideal $I$, $(R/I)\otimes_RM\to M/IM$, given by $\bar r\otimes m\mapsto\overline{rm}$. Here $IM$ consists of finite sums of elements $im$, with $i\in I,m\in M$.
\end{enumerate}
\end{proposition}

\begin{proof}
For (a), multiplication is bilinear, so the displayed map exists. Its inverse sends $m$ to $1\otimes m$. The identity $r\otimes m=1\otimes rm$ checks the composite on a pure tensor.

For (b), exchange the factors by the same construction twice. The composite fixes every pure tensor and therefore every element.

For (c), the displayed family has finite support because $(m_i)_i$ does. It defines a bilinear map and hence a linear map from the tensor product. Conversely, on the summand $M_j\otimes N$ use $\iota_j\otimes\mathrm{id}_N$, where $\iota_j\colon M_j\to\bigoplus_iM_i$ is inclusion. These maps combine into a map from the direct sum. Their composites fix the generating tensors because
\[
(m_i)_i\otimes n=\sum_i\iota_i(m_i)\otimes n.
\]
Only finitely many terms occur; no infinite sum in a module is being taken. The argument includes an empty direct sum, for which both sides are zero.

For (d), changing $r$ by $i\in I$ changes $rm$ by $im\in IM$, so the displayed bilinear map is well defined. In the other direction send $\bar m$ to $\bar1\otimes m$. If $m$ changes by $im'$, its image changes by
$\bar1\otimes im'=\bar i\otimes m'=0$; finite sums give the same conclusion for all of $IM$. The composites fix the generators, using $\bar r\otimes m=\bar1\otimes rm$.

All four constructions commute with module maps: evaluate either order of a proposed naturality square on a pure tensor to obtain the same value. Since pure tensors generate, the square commutes everywhere.
\end{proof}

\begin{example}\label{ex:zmzn}
For positive integers $m,n$,
\[
\Z/m\otimes_\Z\Z/n\cong(\Z/n)/m(\Z/n)
\cong\Z/\gcd(m,n).
\]
The first isomorphism uses (d), and the second uses $m\Z+n\Z=\gcd(m,n)\Z$. Thus $\Z/4\otimes\Z/6\cong\Z/2$, whereas $\Z/2\otimes\Z/3=0$.

More generally, tensoring $\Z^a\oplus\bigoplus_i\Z/m_i$ with $\Z^b\oplus\bigoplus_j\Z/n_j$ gives
\[
\begin{gathered}
\Z^{ab}\oplus\bigoplus_i(\Z/m_i)^b
\oplus\bigoplus_j(\Z/n_j)^a\\
\oplus\bigoplus_{i,j}\Z/\gcd(m_i,n_j).
\end{gathered}
\]
Parts (a)--(c), together with the PID structure theorem of Volume~3, Theorem~7.6, therefore compute every tensor product of finitely generated abelian groups.

For vector spaces with bases $(e_i)_{i\in I}$ and $(f_j)_{j\in J}$, the same direct-sum calculation gives the basis $(e_i\otimes f_j)_{(i,j)\in I\times J}$ of $V\otimes_KW$. In particular $\dim(V\otimes W)=\dim V\dim W$, including zero and infinite cardinal dimensions.
\end{example}

A bilinear map $M\times N\to P$ can also be viewed as a linear map $M\to\operatorname{Hom}(N,P)$: fix $m$ and let $n$ vary. This is the module version of currying in Example~\ref{ex:adjunctions}(c).

\begin{theorem}[tensor and Hom are adjoint]\label{thm:tensorhom}
For modules $M,N,P$, the map
\[
\begin{aligned}
\Phi\colon\operatorname{Hom}(M\otimes_RN,P)
&\longrightarrow\operatorname{Hom}\bigl(M,\operatorname{Hom}(N,P)\bigr),\\
\Phi(u)(m)(n)&=u(m\otimes n)
\end{aligned}
\]
is an $R$-linear isomorphism, natural in $M$ and $P$. Hence $-\otimes_RN$ is left adjoint to $\operatorname{Hom}(N,-)$.
\end{theorem}

\begin{proof}
Separate linearity of $(m,n)\mapsto u(m\otimes n)$ makes $\Phi(u)(m)$ a linear map on $N$, and makes $m\mapsto\Phi(u)(m)$ linear as well. Conversely, from $v\colon M\to\operatorname{Hom}(N,P)$ form the bilinear map $(m,n)\mapsto v(m)(n)$. It induces a linear $\Psi(v)$ with $\Psi(v)(m\otimes n)=v(m)(n)$. These constructions preserve sums and scalars. Substitution gives $\Phi\Psi(v)=v$; the other composite agrees with $u$ on pure tensors, so $\Psi\Phi(u)=u$.

To check naturality, take $\mu\colon M'\to M$ and $\pi\colon P\to P'$. Both routes through the naturality square, evaluated at $m'\in M'$ and $n\in N$, give
\[
\pi\bigl(u(\mu(m')\otimes n)\bigr).
\]
Thus the square commutes, with the variance in $M$ reversed and that in $P$ preserved, as required for the Hom-sets of an adjunction.
\end{proof}

The adjunction gives the exactness statement without a further construction of the tensor product.

\begin{corollary}[tensor is right exact]\label{cor:rightexact}
For an exact sequence $A\xrightarrow f B\xrightarrow g C\to0$ and any module $N$, the sequence
\[
A\otimes N\xrightarrow{f\otimes1}B\otimes N
\xrightarrow{g\otimes1}C\otimes N\to0
\]
is exact.
\end{corollary}

\begin{proof}
Choose an arbitrary target $P$. Part (b) of Proposition~\ref{prop:homleft}, with target $\operatorname{Hom}(N,P)$, gives an exact sequence of Hom-modules. The tensor--Hom adjunction rewrites it as
\[
\begin{gathered}
0\to\operatorname{Hom}(C\otimes N,P)
\xrightarrow{(g\otimes1)^*}\operatorname{Hom}(B\otimes N,P)\\
\xrightarrow{(f\otimes1)^*}\operatorname{Hom}(A\otimes N,P).
\end{gathered}
\]
Naturality is what identifies these arrows with precomposition by $g\otimes1$ and $f\otimes1$; a mere bijection of the three modules would not suffice. Since the displayed sequence is exact for every $P$, Lemma~\ref{lem:testexact} gives the claim.
\end{proof}

The initial injection need not survive. Tensor the opening sequence with $\Z/2$. The map induced by multiplication by $2$ on $\Z$ becomes multiplication by $2$ on $\Z/2$, hence the zero map. The quotient map becomes the identity after the usual identifications. Figure~\ref{fig:v10new-ch03-collapse} displays the resulting exact sequence at the last two terms and its failure at the first. Chapter~7 will measure that failure with $\operatorname{Tor}$.

\begin{figure}[ht]
\centering
\fig{v10-new-ch03-collapse}
\caption{Tensoring with $\Z/2$ preserves the quotient at the right but turns the first injection into zero. The lower row is exact at its middle and right non-zero terms; it is not exact at its first $\Z/2$.}
\label{fig:v10new-ch03-collapse}
\end{figure}


% V10-EXPANSION WORKED EXAMPLE 2
\subsection*{Worked example: tensoring a module given by relations}

Let
\[
M=\Z^2/\langle(2,0),(1,3)\rangle,
\]
and write $a,b$ for the classes of the coordinate vectors. Its relations are $2a=0$ and $a+3b=0$. The second expresses $a=-3b$, and substitution into the first gives $6b=0$. Thus $b$ generates $M$, and its order divides six.

There is no additional hidden relation. Send $(r,s)\in\Z^2$ to $3r+s$ modulo $6$. Both defining relation vectors go to zero, so this induces $M\to\Z/6$ sending $b$ to $\bar1$. Conversely, $\bar1\mapsto b$ defines $\Z/6\to M$ because $6b=0$. The two maps fix the generators and are inverse. Hence $M\cong\Z/6$.

Now tensor with $N=\Z/4$. Right exactness tells us how to use the presentation directly:
\[
(\Z/4)^2\xrightarrow{\,D\,}(\Z/4)^2
 \longrightarrow M\otimes_\Z\Z/4\longrightarrow0,
\qquad D=\begin{pmatrix}2&1\\0&3\end{pmatrix}.
\]
The columns of $D$ are the original relations, now read modulo four. The tensor product is the cokernel of this map, not its kernel. In that cokernel, write $a',b'$ for the two coordinate classes. Again $a'=-3b'$, while $2a'=0$ gives $6b'=0$. Together with $4b'=0$, this gives $2b'=0$. The map $(r,s)\mapsto r+s$ modulo $2$ kills the two columns of $D$ and sends $b'$ to $\bar1$. It proves that $b'$ is not zero. Therefore the cokernel is exactly $\Z/2$, agreeing with
\[
(\Z/6)\otimes_\Z(\Z/4)\cong\Z/\gcd(6,4).
\]
This is useful even when a presentation has many generators: tensor the free modules, keep the relation matrix, and compute its cokernel over the new coefficients. The first map may gain a kernel. Right exactness makes no promise that a relation matrix remains injective; it promises that the quotient it presents is still the correct quotient.

\section{Projective, injective and flat modules}

The example asks which modules avoid these losses. There are three useful answers, depending on whether we use covariant Hom, contravariant Hom, or tensor.

\begin{definition}\label{def:projective}
A module $P$ is \emph{projective} when $\operatorname{Hom}(P,-)$ preserves short exact sequences. By left exactness, this says precisely that every map $P\to C$ lifts to $P\to B$ along every surjection $B\to C$.

A module $E$ is \emph{injective} when $\operatorname{Hom}(-,E)$ preserves short exact sequences, with the arrows reversed. Equivalently, every map $A\to E$ extends to $B\to E$ along every injection $A\to B$.

A module $N$ is \emph{flat} when $-\otimes_RN$ preserves short exact sequences. Since tensor is right exact, this is equivalent to requiring that $f\otimes1_N$ be injective whenever $f$ is injective.
\end{definition}

A lift and an extension are different requests. A lift chooses preimages through a surjection; an extension supplies values on a larger domain. Figure~\ref{fig:v10new-ch03-lifts} puts the two requests beside each other. The dashed arrow is the one the definition promises, and it need not be unique.

\begin{figure}[ht]
\centering
\fig{v10-new-ch03-lifts}
\caption{Projectivity supplies a lift $P\to B$ through a surjection $B\to C$. Injectivity supplies an extension $B\to E$ from a submodule $A\subseteq B$. In both pictures the two routes from the upper left to the lower right must agree.}
\label{fig:v10new-ch03-lifts}
\end{figure}

\begin{proposition}\label{prop:projective}
\leavevmode
\begin{enumerate}[label=(\alph*)]
\item Every free module is projective. A module is projective exactly when it is a direct summand of a free module.
\item Every free module is flat, and every direct summand of a flat module is flat. In particular every projective module is flat.
\item If $R$ is a domain, every flat module is torsion-free: $rx=0$, $r\ne0$, forces $x=0$.
\end{enumerate}
\end{proposition}

\begin{proof}
\emph{(a) Lift a basis, then split a free presentation.} Let $F$ have basis $(e_i)$ and let $g\colon B\to C$ be onto. Given $u\colon F\to C$, choose $b_i$ with $g(b_i)=u(e_i)$. The assignments $e_i\mapsto b_i$ extend uniquely to a linear map $\tilde u$, and $g\tilde u=u$ because both maps agree on the basis. This includes the empty basis of the zero module. For an infinite basis the choices use the axiom of choice, as did the basis constructions of Volume~7.

For any module $P$, a free module with one basis vector for each element of $P$ surjects onto $P$. If $P$ is projective, lift $\mathrm{id}_P$ through that surjection $F\to P$. The splitting lemma gives $F\cong\ker(F\to P)\oplus P$.

Conversely, suppose $F=P\oplus Q$ is free. Given $u\colon P\to C$ and an onto $B\to C$, extend $u$ to $F$ by zero on $Q$, lift the extension using freeness, and restrict that lift to $P$. This proves projectivity.

\emph{(b) Tensor works one direct summand at a time.} Part (c) of Proposition~\ref{prop:tensorrules}, and the symmetry in part (b), identify $f\otimes1_{N\oplus N'}$ with the direct sum of $f\otimes1_N$ and $f\otimes1_{N'}$. Such a sum is injective exactly when each summand map is injective. The same assertion holds for an arbitrary direct sum because every element has finite support. Since $-\otimes R$ is naturally the identity, $R$ is flat, all direct sums of copies of $R$ are flat, and summands of flat modules are flat. Combine this with (a).

\emph{(c) Test multiplication by a non-zero scalar.} In a domain the map $R\xrightarrow{\,r\,}R$ is injective for $r\ne0$. Tensoring it with a flat $N$ remains injective. Under $R\otimes N\cong N$, it is the map $x\mapsto rx$. Its kernel is therefore zero.
\end{proof}

Over a ring with zero divisors, the domain hypothesis in (c) cannot be dropped: even the free module $R$ can contain a non-zero $x$ with $rx=0$ and $r\ne0$. Projectivity, by contrast, is a lifting property over any ring in our standing convention.

To test injectivity it is enough to extend maps from ideals. Baer's criterion replaces an arbitrary submodule inclusion by this much smaller family of tests.

\begin{theorem}[Baer's criterion]\label{thm:baer}
An $R$-module $E$ is injective if and only if each linear map $I\to E$ from an ideal $I\subseteq R$ extends to a linear map $R\to E$.
\end{theorem}

\begin{proof}
An injective module has the stated property by applying its definition to $I\subseteq R$.

Conversely, assume the ideal-extension property. Given $A\subseteq B$ and $u\colon A\to E$, consider pairs $(A_1,u_1)$ with $A\subseteq A_1\subseteq B$ and $u_1\colon A_1\to E$ extending $u$. Order them by extension. The pair $(A,u)$ makes the set nonempty. For a chain, the union of the domains is a submodule, and the maps agree on overlaps; their union is a linear extension on that union. The empty chain has $(A,u)$ as an upper bound. Zorn's lemma, Volume~7, Theorem~1.19, therefore supplies a maximal pair $(A_1,u_1)$.

If $b\in B\setminus A_1$, set
\[
I=\{r\in R:rb\in A_1\}.
\]
This is an ideal: differences and multiples still send $b$ into $A_1$. The map $r\mapsto u_1(rb)$ is linear on $I$. By hypothesis it extends to $w\colon R\to E$. Define
\[
u_2(a+rb)=u_1(a)+w(r),\qquad a\in A_1, r\in R.
\]
To check this formula, suppose $a+rb=a'+r'b$. Then $r-r'\in I$, and
\[
w(r-r')=u_1((r-r')b)=u_1(a'-a).
\]
Rearranging gives $u_1(a)+w(r)=u_1(a')+w(r')$. Thus $u_2$ is well defined. It is linear by the formula, extends $u_1$, and has the strictly larger domain $A_1+Rb$. This contradicts maximality. Hence $A_1=B$, so every $u$ extends and $E$ is injective.
\end{proof}

\begin{corollary}\label{cor:divisible}
Over a principal ideal domain, an $R$-module $E$ is injective exactly when it is \emph{divisible}: for each $e\in E$ and each non-zero $r\in R$, an $e'\in E$ satisfies $re'=e$. In particular, $\Q$ and $\Q/\Z$ are injective as abelian groups, whereas $\Z$ is not.
\end{corollary}

\begin{proof}
Each non-zero ideal has the form $(r)$ with $r\ne0$. Since $R$ is a domain, a map $u\colon(r)\to E$ is determined by the arbitrary value $e=u(r)$: define $u(sr)=se$. Extending it to $R$ asks for a value $e'=w(1)$ with $re'=e$. Thus all such maps extend exactly when $E$ is divisible. The only map from the zero ideal also extends, by zero. Baer's criterion proves the equivalence.

Every rational number can be divided by a non-zero integer, so $\Q$ is divisible. A quotient of a divisible group is divisible: divide a representative and then take its class. This applies to $\Q/\Z$. The element $1\in\Z$ cannot be divided by $2$ within $\Z$, so $\Z$ is not divisible and not injective.
\end{proof}

\begin{theorem}\label{thm:pidflat}
For a finitely generated module $M$ over a principal ideal domain, the following four conditions are equivalent: $M$ is free, $M$ is projective, $M$ is flat, and $M$ is torsion-free.
\end{theorem}

\begin{proof}
Proposition~\ref{prop:projective} gives the implications from free to projective to flat to torsion-free. For the last implication back to free, use the PID structure theorem from Volume~3, Theorem~7.6:
\[
M\cong R^n\oplus R/(d_1)\oplus\cdots\oplus R/(d_k),
\]
where each $d_i$ is a non-zero non-unit. The non-zero class of $1$ in $R/(d_i)$ is killed by $d_i$. A torsion-free $M$ can therefore have no such summand. What remains is $R^n$, including $R^0=0$.
\end{proof}

Finite generation is essential. The group $\Q$ is torsion-free but is neither free nor projective, as Exercise~\ref{ex:qnotproj} shows. It is also flat, a fact that can already be checked here. For any abelian group $A$, identify
\[
A\otimes_\Z\Q\cong S^{-1}A,\qquad S=\{1,2,3,\ldots\},
\]
where $S^{-1}A$ consists of fractions $a/s$, with $a/s=a'/s'$ when some $t\ge1$ kills $s'a-sa'$. Addition uses common denominators, and the maps in the identification are
\[
a\otimes(r/s)\longmapsto ra/s,\qquad
 a/s\longmapsto a\otimes(1/s).
\]
The first respects bilinearity. For the second, the equality condition gives
$(s'a-sa')\otimes1/(ss')=0$, because a tensor with a finite-order element and a divisible rational factor is zero: if $tx=0$, then $x\otimes q=tx\otimes(q/t)=0$. The two maps are inverse on generators. If $A\to B$ is injective and $a/s$ maps to zero, some $t\ge1$ kills the image of $a$, hence kills $a$ itself. Thus $a/s=0$. Tensoring an injection with $\Q$ remains injective, proving flatness without waiting for the localisation chapter.

For finite-dimensional vector spaces $V,W$, there is another useful identification. The bilinear map
\[
V^*\times W\to\operatorname{Hom}(V,W),\qquad
(\varphi,w)\mapsto(v\mapsto\varphi(v)w)
\]
gives a linear $V^*\otimes W\to\operatorname{Hom}(V,W)$. On bases, $e_i^*\otimes f_j$ maps to the elementary map sending $e_i$ to $f_j$ and the other basis vectors to zero. These elementary maps form a basis, so the map is an isomorphism. A pure tensor corresponds to a map of rank at most $1$. A sum of $r$ pure tensors therefore has rank at most $r$; conversely, a rank-$r$ map can be written using a basis of its image and $r$ coordinate functionals. Hence matrix rank is exactly two-factor tensor rank. For three factors this argument no longer reduces the problem to the image of a single linear map.

\begin{beauty}
Matrix multiplication is bilinear, so it determines a linear map $M_2(K)\otimes M_2(K)\to M_2(K)$. Because these spaces are finite-dimensional, that map is the same data as a tensor in $M_2(K)^*\otimes M_2(K)^*\otimes M_2(K)$. Strassen's seven products in Theorem~\ref{linear:thm:strassen} express this tensor as a sum of seven pure tensors. Winograd's lower bound, quoted in Chapter~\ref{linear:ch:tensors}, says that fewer do not suffice. The universal property explains why changing bases changes the formulas without changing the bilinear problem they solve.
\end{beauty}

\section{What we proved, and what comes next}

An exact sequence specifies which elements disappear and where they came from. Choosing compatible lifts is a separate question: it is the difference between the split sequence $A\to A\oplus C\to C$ and the sequence $\Z\xrightarrow2\Z\to\Z/2$.

The snake construction gives a useful way to handle the choices that cannot be made uniquely. Lift an element, move it through a commutative square, and take its class in a quotient. The quotient removes exactly the ambiguity in the lift. That is why the connecting homomorphism is well defined.

Hom retains kernels, and tensor retains cokernels and surjections. Projective and injective modules let us repair the missing Hom surjections by lifting or extending maps. Flat modules let us keep injections after tensoring. Chapter~7 will record the failures for arbitrary modules as $\operatorname{Tor}$ and $\operatorname{Ext}$.

Before measuring those failures, we need to know more about the rings on which they occur. Can infinitely many polynomial equations really impose more conditions than a finite list? The next chapter answers this with Hilbert's basis theorem.

\section{Exercises}

\begin{exercise}[$\star$]\label{ex:zmzncomp}
For positive integers $m,n$, compute $\operatorname{Hom}_\Z(\Z/m,\Z/n)$ as an abelian group and prove that it is isomorphic to $\Z/\gcd(m,n)$. Show also that $\Q\otimes_\Z(\Q/\Z)=0$. More generally, prove $D\otimes_\Z T=0$ whenever $D$ is divisible and $T$ is torsion, meaning that each of its elements is killed by some non-zero integer.
\end{exercise}

\begin{exercise}[$\star$]\label{ex:whichsplit}
Decide whether each short exact sequence splits, and justify the answer:
\[
\begin{gathered}
0\to\Z/2\to\Z/6\to\Z/3\to0,\\
0\to\Z/3\to\Z/9\to\Z/3\to0,\\
0\to\Z^2\xrightarrow{u}\Z^2\to\Z/5\to0,
\qquad u(x,y)=(x,5y),\\
0\to K\to\Z^2\xrightarrow{v}\Z\to0.
\end{gathered}
\]
In the first two sequences use the subgroup inclusions and residue quotient maps; in the third the quotient is $(x,y)\mapsto\bar y$. In the last, $v$ is any surjection and $K=\ker v$.
\end{exercise}

\begin{exercise}[$\star\star$]\label{ex:composite}
For homomorphisms $f\colon A\to B$ and $g\colon B\to C$, construct the arrows in
\[
\begin{gathered}
0\to\ker f\to\ker(gf)\to\ker g\to\operatorname{coker}f\\
\to\operatorname{coker}(gf)\to\operatorname{coker}g\to0
\end{gathered}
\]
and prove exactness at every term.
\end{exercise}

\begin{exercise}[$\star\star$]\label{ex:qnotproj}
\leavevmode
\begin{enumerate}[label=(\alph*)]
\item Prove $\operatorname{Hom}_\Z(\Q,\Z)=0$. Use the direct-summand characterisation of projective modules to show that $\Q$ is not projective.
\item Prove that any finite direct sum of injective modules is injective. Establish the decomposition
\[
\Q/\Z\cong\bigoplus_{p\ \mathrm{prime}}\Z[1/p]/\Z.
\]
\end{enumerate}
\end{exercise}

\begin{exercise}[$\star\star$]\label{ex:tensoralgebra}
For commutative $R$-algebras $A,B$, use the multiplication $(a\otimes b)(a'\otimes b')=aa'\otimes bb'$ on $A\otimes_RB$; you may assume that it makes this a commutative ring. Prove the ring isomorphisms
\[
\C\otimes_\R\C\cong\C\times\C,
\qquad
\Q(\sqrt2)\otimes_\Q\Q(\sqrt3)\cong\Q(\sqrt2,\sqrt3).
\]
In particular, explain why the second tensor product is a field. Derive and use $R[X]/(p)\otimes_RA\cong A[X]/(p)$.
\end{exercise}

\begin{exercise}[$\star\star$]\label{ex:fivelemma}
In a commutative diagram with exact rows $A_1\to A_2\to A_3\to A_4\to A_5$ and $B_1\to B_2\to B_3\to B_4\to B_5$, write $h_i\colon A_i\to B_i$ for the vertical maps. Suppose $h_2,h_4$ are isomorphisms, $h_1$ is onto, and $h_5$ is injective. Prove the \emph{five lemma}: $h_3$ is an isomorphism.
\end{exercise}

\begin{exercise}[$\star\star$]\label{ex:v10new-ch03-quotients}
Let $I,J$ be ideals of $R$. Construct a natural isomorphism
\[
(R/I)\otimes_R(R/J)\cong R/(I+J)
\]
by specifying it on pure tensors and writing its inverse. Compute $(\Z/12)\otimes_\Z(\Z/18)$. Then take $R=K[\varepsilon]/(\varepsilon^2)$ and compute $(R/(\varepsilon))\otimes_R(R/(\varepsilon))$. Explain why the last tensor product does not recover the non-zero nilpotent element of $R$.
\end{exercise}

\begin{exercise}[$\star\star$]\label{ex:v10new-ch03-snake}
Let $m,n\ge1$, put $d=\gcd(m,n)$, and form a diagram with both rows $0\to\Z\xrightarrow{m}\Z\to\Z/m\to0$ and all three vertical maps multiplication by $n$. Compute the snake map
\[
\partial\colon\ker(n\colon\Z/m\to\Z/m)\longrightarrow\Z/n.
\]
Show that it sends the class of $m/d$ to the class of $n/d$, that it is injective, and that its image is the kernel of multiplication by $m$ on $\Z/n$. Work out $m=6,n=4$ explicitly.
\end{exercise}

\begin{exercise}[$\star\star$]\label{ex:v10new-ch03-idempotent}
An element $e\in R$ is \emph{idempotent} if $e^2=e$. Show that $R=eR\oplus(1-e)R$ as modules, and hence that $eR$ is projective and flat. Construct an isomorphism $eR\otimes_RM\cong eM$ for every module $M$. For $R=\Z/6$ and $e=\bar3$, identify $eR$ and $(1-e)R$ as abelian groups. Why does the torsion in $eR$ not contradict Proposition~\ref{prop:projective}(c)?
\end{exercise}

\begin{exercise}[$\star\star\star$]\label{ex:v10new-ch03-finiterank}
For arbitrary vector spaces $V,W$ over $K$, prove that
\[
V^*\otimes_KW\longrightarrow\operatorname{Hom}_K(V,W),
\qquad \varphi\otimes w\longmapsto(v\mapsto\varphi(v)w)
\]
is injective and has image exactly the maps of finite rank. For $V=W=\bigoplus_{j\ge1}Ke_j$, show that the identity is not in the image. Explain why the map is surjective whenever either $V$ or $W$ is finite-dimensional, including the zero-dimensional cases.
\end{exercise}


\begin{exercise}[$\star$]\label{ex:v10exp-ch03-presentation}
Let $M=\Z^2/\langle(4,0),(2,6)\rangle$. Find its cyclic direct-sum decomposition by explicit integer row and column operations on its relation matrix. Compute $M\otimes_\Z\Z/6$ and $\operatorname{Hom}_\Z(M,\Z/6)$ as abelian groups, and count the homomorphisms in the latter group.
\end{exercise}

\begin{exercise}[$\star\star$]\label{ex:v10exp-ch03-projective-splitting}
Prove that an $R$-module $P$ is projective if and only if every short exact sequence ending in $P$ splits. For the reverse direction, start from a surjection $q:B\to C$ and a map $u:P\to C$, and consider $E=\{(p,b)\in P\oplus B:u(p)=q(b)\}$. Give the maps explicitly rather than assuming a lifting property for $E$.
\end{exercise}

\begin{exercise}[$\star\star$]\label{ex:v10exp-ch03-pid-hypothesis}
For a field $K$, let $R=K[X,Y]$ and $I=(X,Y)\subset R$. Prove that $I$ is a finitely generated torsion-free $R$-module but is not free. Explain why this does not contradict Theorem~\ref{thm:pidflat}. You may use the earlier UFD theorem for polynomial rings, but no local-ring theorem from Chapter~5 is needed.
\end{exercise}

\begin{exercise}[$\star\star\star$]\label{ex:v10exp-ch03-extensions}
For positive integers $n,k$ and an integer $t$, let the abelian group $E_t$ have generators $a,b$ and relations $na=0$, $kb=ta$. Show that $a$ has order exactly $n$ and that the maps $\Z/n\to E_t$, $\bar1\mapsto a$, and $E_t\to\Z/k$, $a\mapsto0$, $b\mapsto\bar1$, form a short exact sequence. Determine exactly when it splits and count all its sections. Compare $n=k=2$ for $t=0$ and $t=1$.
\end{exercise}

\section{Hints}

\begin{hint}{ex:zmzncomp}
Choose the image of $\bar1$ and impose the relation $m\bar1=0$. For a pure tensor $d\otimes t$, first choose a non-zero integer $n$ killing $t$, then divide $d$ by $n$ in $D$.
\end{hint}

\begin{hint}{ex:whichsplit}
A section of the first quotient sends $\bar1$ to $\bar4$. For the next two sequences compare torsion or element orders in the proposed direct sums. For the last, lift $1\in\Z$ and extend linearly.
\end{hint}

\begin{hint}{ex:composite}
The arrows after $\ker(gf)$ are induced by $f$, by the quotient $B\to B/f(A)$, by $g$, and by the quotient $C/gf(A)\to C/g(B)$. Check that the map induced by $g$ is well defined. The snake lemma also applies to the rows $A\to B\to\operatorname{coker}f\to0$ and $0\to C\xrightarrow{\mathrm{id}}C\to0$; its connecting arrow starts at $\operatorname{coker}f$.
\end{hint}

\begin{hint}{ex:qnotproj}
The integer image of a rational number would be divisible by every positive integer. For the decomposition, split a denominator into relatively prime prime powers and apply Bézout's identity.
\end{hint}

\begin{hint}{ex:tensoralgebra}
First tensor the presentation of $R[X]/(p)$ and use right exactness. Over $\C$, factor $X^2+1$ and apply the Chinese remainder theorem. Over $\Q(\sqrt2)$, test whether $(a+b\sqrt2)^2$ can equal $3$ for rational $a,b$.
\end{hint}

\begin{hint}{ex:fivelemma}
For injectivity, send $x\in\ker h_3$ to $A_4$, then lift it back through $A_2$. For surjectivity, send $y\in B_3$ to $B_4$, lift through $h_4$, and correct the remaining error through $B_2$.
\end{hint}

\begin{hint}{ex:v10new-ch03-quotients}
The map should send $\bar a\otimes\bar b$ to $\overline{ab}$. The proposed inverse sends $\bar r$ to $\bar1\otimes\bar r$. Check separately what happens to a summand from $I$ and one from $J$.
\end{hint}

\begin{hint}{ex:v10new-ch03-snake}
If $\bar k$ lies in the kernel, $m$ divides $nk$. Lift it by the integer $k$, multiply by $n$, and pull back through multiplication by $m$. The final answer is a class modulo $n$.
\end{hint}

\begin{hint}{ex:v10new-ch03-idempotent}
Write $r=er+(1-e)r$ and multiply an element of the intersection by $e$. The tensor map is $er\otimes m\mapsto erm$; its inverse on $eM$ sends $em$ to $e\otimes em$.
\end{hint}

\begin{hint}{ex:v10new-ch03-finiterank}
A tensor uses only finitely many vectors of $W$. Choose a basis of their span and rewrite the tensor using that basis. Independent output vectors force the corresponding input functionals to vanish separately.
\end{hint}


\begin{hint}{ex:v10exp-ch03-presentation}
The relation matrix has columns $(4,0)$ and $(2,6)$. Swap them, subtract three times the first row from the second, then clear the remaining upper-right entry.
\end{hint}
\begin{hint}{ex:v10exp-ch03-projective-splitting}
The projection $E\to P$ is onto because $q$ is onto. A section has second coordinate equal to the lift being sought.
\end{hint}
\begin{hint}{ex:v10exp-ch03-pid-hypothesis}
Pass to the fraction field to show that a free basis of $I$ would have exactly one element. A principal generator would divide both $X$ and $Y$ in a UFD.
\end{hint}
\begin{hint}{ex:v10exp-ch03-extensions}
Use the quotient of $\Z^2$ by relation vectors $(n,0)$ and $(-t,k)$. Every lift of $\bar1\in\Z/k$ has the form $b+sa$, and it defines a section precisely when $t+ks$ vanishes modulo $n$.
\end{hint}

\section{Solutions}

\begin{solution}{ex:zmzncomp}
\idea A map out of a cyclic group is determined by one element. In a tensor, divisibility lets us move an annihilating integer onto the torsion factor.
\writeup Put $d=\gcd(m,n)$. A homomorphism $\Z/m\to\Z/n$ is determined by $x=u(\bar1)$ with $mx=0$. The solutions are precisely
\[
x=\overline{j(n/d)},\qquad j=0,\ldots,d-1.
\]
Indeed, writing $m=dm_0,n=dn_0$ with $(m_0,n_0)=1$, the congruence $n\mid mx$ is equivalent to $n_0\mid x$. Addition of homomorphisms adds these images, so evaluation at $\bar1$ identifies the Hom-group with the cyclic subgroup generated by $\overline{n/d}$, of order $d$.

For $d_0\in D$ and $t\in T$, choose $n\ge1$ with $nt=0$, and choose $d_1\in D$ with $nd_1=d_0$. Then
\[
d_0\otimes t=nd_1\otimes t=d_1\otimes nt=0.
\]
Every element is a sum of pure tensors, so the whole tensor product is zero. The group $\Q$ is divisible and $\Q/\Z$ is torsion, giving the requested instance.
\end{solution}

\begin{solution}{ex:whichsplit}
\idea Use a section when a lift of the quotient generator has the correct order. Otherwise compare the middle term with the direct sum a splitting would require.
\writeup In the first sequence $s(\bar1)=\bar4\in\Z/6$ defines a section: $3\bar4=0$ and $\bar4$ maps to $\bar1$ modulo $3$. Thus it splits.

The second would require $\Z/9\cong\Z/3\oplus\Z/3$. The right side is killed by $3$, but the class of $1$ on the left is not. Thus it does not split. The third would require $\Z^2\cong\Z^2\oplus\Z/5$, which is impossible because the right side has non-zero torsion and the left side is torsion-free.

For the last sequence, choose $b\in\Z^2$ with $v(b)=1$. The map $s(k)=kb$ is a section, so it splits. The kernel is a free rank-one group: if $v(x,y)=ax+by$, surjectivity says $(a,b)=1$, and $K$ is generated by $(-b,a)$. This also checks directly that the middle group has the required direct-sum form.
\end{solution}

\begin{solution}{ex:composite}
\idea Quotient maps provide the arrows on the right. Exactness amounts to lifting a representative and subtracting something from the preceding image.
\writeup Write the sequence as
\[
\begin{gathered}
0\to\ker f\xrightarrow{\iota}\ker(gf)\xrightarrow f\ker g
\xrightarrow{q}B/f(A)\\
\xrightarrow{\bar g}C/gf(A)\xrightarrow{p}C/g(B)\to0,
\end{gathered}
\]
where $q(b)=b+f(A)$, $\bar g(b+f(A))=g(b)+gf(A)$, and $p$ is the natural quotient. Changing $b$ by $f(a)$ changes $g(b)$ by $gf(a)$, so $\bar g$ is well defined. The other maps are restrictions or quotient maps and are linear.

The map $\iota$ is injective. An $a\in\ker(gf)$ maps to zero under $f$ exactly when $a\in\ker f$. An element $b\in\ker g$ is killed by $q$ exactly when $b=f(a)$; then $gf(a)=0$, so $a\in\ker(gf)$.

If $\bar g([b])=0$, write $g(b)=gf(a)$. Then $b-f(a)\in\ker g$ and $q(b-f(a))=[b]$. Conversely, $g$ kills every element of $\ker g$, so $\bar gq=0$. This proves exactness at $B/f(A)$.

A class $[c]\in C/gf(A)$ lies in $\ker p$ precisely when $c\in g(B)$, which says that it is in $\operatorname{im}\bar g$. Finally $p$ is onto because every class modulo $g(B)$ has a representative in $C$. These checks prove exactness everywhere.

In the snake-lemma derivation from the hint, the connecting map is $\bar g\colon\operatorname{coker}f\to\operatorname{coker}(gf)$. The map $\ker g\to\operatorname{coker}f$ precedes it and is induced by the quotient of $B$.
\end{solution}

\begin{solution}{ex:qnotproj}
\idea Coordinates of a free abelian group take integer values. A divisible source cannot have a non-zero image in any such coordinate.
\writeup (a) If $u\colon\Q\to\Z$, then for $q\in\Q$ and every $n\ge1$,
\[
u(q)=n\,u(q/n).
\]
An integer divisible by every $n$ is zero: for a non-zero integer $a$, choose $n>|a|$. Thus $u=0$. If $\Q$ were projective, it would embed as a summand in some free abelian group $\bigoplus_{s\in S}\Z$. A non-zero image vector has a non-zero coordinate; composing the embedding with that coordinate projection would give a non-zero map $\Q\to\Z$, a contradiction. Since free modules are projective, this also proves that $\Q$ is not free.

(b) For finitely many injective modules $E_1,\ldots,E_r$, extend each component of a map $A\to\bigoplus_iE_i$ across $A\subseteq B$. The finitely many extensions combine to a map into that direct sum. The empty sum is zero and is injective as well.

For the decomposition of $\Q/\Z$, each summand $\Z[1/p]/\Z$ embeds through its rational representatives. Given a fraction $a/n$, factor $n=\prod_{i=1}^r p_i^{k_i}$. The integers $n/p_i^{k_i}$ have greatest common divisor $1$, so choose $c_i$ with
\[
\sum_i c_i\frac{n}{p_i^{k_i}}=1.
\]
Then $a/n=\sum_i ac_i/p_i^{k_i}$ is a finite sum of classes from the indicated summands. Thus the sum map is onto.

For uniqueness, suppose $\sum_{i=1}^r x_i=0$ where the order of $x_i$ is a power of the distinct prime $p_i$. For a fixed $i$, multiplying by a product of sufficiently large powers of the other primes kills all $x_j$, $j\ne i$. That product is relatively prime to the order of $x_i$, so multiplication by it is invertible on the cyclic group generated by $x_i$, and $x_i=0$. The sum map is injective. Each summand is divisible: multiplication by a prime different from $p$ is invertible on each finite $p$-power subgroup, and division by $p$ is obtained by increasing a denominator. Hence each is injective by Corollary~\ref{cor:divisible}.
\end{solution}

\begin{solution}{ex:tensoralgebra}
\idea A polynomial presentation turns tensoring into a change of coefficient ring. Factorisation of the defining polynomial then decides whether the result splits or remains a field.
\writeup As an $R$-module, $R[X]=\bigoplus_{j\ge0}RX^j$, so Proposition~\ref{prop:tensorrules} gives $R[X]\otimes_RA\cong A[X]$. Tensor the right-exact presentation
\[
R[X]\xrightarrow{\,p\,}R[X]\to R[X]/(p)\to0.
\]
The induced map on $A[X]$ is multiplication by $p$, now with coefficients mapped to $A$. Its cokernel is $A[X]/(p)$. The resulting isomorphism sends $(f+(p))\otimes a$ to $af+(p)$; checking products of pure tensors shows that it is a ring isomorphism.

Apply this to $\C=\R[X]/(X^2+1)$. Over $\C$, the factors $X-i$ and $X+i$ generate comaximal ideals because their difference is the non-zero constant $2i$. The Chinese remainder theorem, Volume~3, Theorem~6.9, gives
\[
\C\otimes_\R\C\cong\C[X]/(X^2+1)\cong\C\times\C.
\]

For the second example, use $\Q(\sqrt3)=\Q[X]/(X^2-3)$ and change coefficients to $\Q(\sqrt2)$. If $a+b\sqrt2$, with rational $a,b$, squared to $3$, comparison of coefficients would give
\[
a^2+2b^2=3,\qquad 2ab=0.
\]
If $b=0$, a rational square would be $3$, impossible by its odd $3$-adic exponent. If $a=0$, a rational square would be $3/2$, again impossible by its prime exponents. Thus $X^2-3$ has no root over $\Q(\sqrt2)$ and, being quadratic, is irreducible there. Its quotient is the field $\Q(\sqrt2,\sqrt3)$.
\end{solution}

\begin{solution}{ex:fivelemma}
\idea Prove injectivity and surjectivity separately. Exactness moves a vanishing image one position left; the hypotheses on the vertical maps let us move between rows.
\writeup Denote the horizontal arrows by $\alpha_i\colon A_i\to A_{i+1}$ and $\beta_i\colon B_i\to B_{i+1}$. The squares say $h_{i+1}\alpha_i=\beta_i h_i$.

For injectivity, let $h_3(x)=0$. Then $h_4(\alpha_3x)=\beta_3h_3x=0$, and injectivity of $h_4$ gives $\alpha_3x=0$. Write $x=\alpha_2w$. Now $\beta_2h_2w=h_3x=0$, so $h_2w=\beta_1v'$ for some $v'\in B_1$. Choose $v\in A_1$ with $h_1v=v'$. Commutativity gives
\[
h_2(\alpha_1v)=\beta_1h_1v=\beta_1v'=h_2w.
\]
Since $h_2$ is injective, $w=\alpha_1v$, and $x=\alpha_2\alpha_1v=0$.

For surjectivity, take $y\in B_3$. Choose $z\in A_4$ with $h_4z=\beta_3y$. Then
\[
h_5(\alpha_4z)=\beta_4h_4z=\beta_4\beta_3y=0.
\]
Injectivity of $h_5$ gives $\alpha_4z=0$, so $z=\alpha_3x$ for some $x\in A_3$. Hence $\beta_3(y-h_3x)=0$. Write $y-h_3x=\beta_2u'$ with $u'\in B_2$, and lift $u'=h_2u$. Finally
\[
y=h_3x+\beta_2h_2u=h_3(x+\alpha_2u).
\]
Thus $h_3$ is onto and injective.
\end{solution}

\begin{solution}{ex:v10new-ch03-quotients}
\idea Both ideals must act by zero on the tensor product. Their sum is exactly the relation needed on the resulting cyclic module.
\writeup The rule $\bar a\otimes\bar b\mapsto ab+(I+J)$ is bilinear and independent of representatives: changing $a$ by an element of $I$, or $b$ by an element of $J$, changes the product by an element of $I+J$. It therefore defines a linear map $u$ to $R/(I+J)$.

Define $v(r+(I+J))=\bar1\otimes\bar r$. If $i\in I$, then $\bar1\otimes\bar i=\bar i\otimes\bar1=0$ because $i$ is zero in the first factor. A $j\in J$ is already zero in the second factor. Thus $v$ is well defined and linear. The composite $uv$ fixes $\bar r$. Conversely,
\[
vu(\bar a\otimes\bar b)=\bar1\otimes\overline{ab}
=\bar a\otimes\bar b.
\]
The maps are inverse, and their formulas commute with all maps induced by ring maps respecting the ideals.

For integers $(12)+(18)=(6)$, so the tensor product is $\Z/6$. For the ring with $\varepsilon^2=0$ and $I=J=(\varepsilon)$, the answer is $R/(\varepsilon)\cong K$. The tensor factors already make $\varepsilon$ act as zero; taking their tensor product cannot restore its non-zero class in $R$.
\end{solution}

\begin{solution}{ex:v10new-ch03-snake}
\idea Follow an integer representative through the snake construction, then express the kernel as a cyclic subgroup.
\writeup For $\bar k\in\ker n$ on $\Z/m$, the integer $nk$ is divisible by $m$. Lift $\bar k$ to $k\in\Z$, map it down to $nk$, and pull it back through multiplication by $m$. This gives
\[
\partial(\bar k)=\overline{nk/m}\quad\hbox{in }\Z/n.
\]
Changing $k$ to $k+mt$ changes this value by $nt$, so the formula is well defined. Write $m=dm_0,n=dn_0$ with $(m_0,n_0)=1$. The condition $m\mid nk$ is equivalent to $m_0\mid k$. Thus the kernel is generated by the class of $m_0=m/d$, of order $d$. Its image is the class of $n_0=n/d$, also of order $d$, so $\partial$ is injective. The elements killed by $m$ in $\Z/n$ are exactly the multiples of $n_0$, proving the image assertion.

For $m=6,n=4$, the kernel in $\Z/6$ is $\{\bar0,\bar3\}$, and the image in $\Z/4$ is $\{\bar0,\bar2\}$. The map sends $\bar3$ to $\bar2$. If $d=1$, both groups are zero; the same formula still applies.
\end{solution}

\begin{solution}{ex:v10new-ch03-idempotent}
\idea An idempotent separates a module into the part on which it acts as identity and the part on which it acts as zero.
\writeup Every $r$ equals $er+(1-e)r$. If $x\in eR\cap(1-e)R$, then $ex=x$ from the first membership, while $ex=e(1-e)r=0$ from the second. Thus the sum is direct. The summand $eR$ of the free module $R$ is projective and flat by Proposition~\ref{prop:projective}.

Multiplication gives $u\colon eR\otimes_RM\to eM$, $er\otimes m\mapsto erm$. Define $v\colon eM\to eR\otimes_RM$ by $v(x)=e\otimes x$ for $x\in eM$. This is a well-defined linear map on the submodule $eM$ itself. Its composite with $u$ sends $x$ to $ex=x$. On a pure tensor,
\[
vu(er\otimes m)=e\otimes erm=er\otimes em=er\otimes m,
\]
since $er\otimes em=e(er)\otimes m=er\otimes m$. Hence $u,v$ are inverse.

Modulo $6$, $e=3$ satisfies $e^2=e$, and $eR=\{0,3\}\cong\Z/2$ as an abelian group. The other summand is $4R=\{0,2,4\}\cong\Z/3$. Proposition~\ref{prop:projective}(c) assumes a domain; $\Z/6$ has non-zero zero divisors. Also projectivity here is over $\Z/6$, not over $\Z$.
\end{solution}

\begin{solution}{ex:v10new-ch03-finiterank}
\idea A finite tensor expression has a finite-dimensional output space. Conversely, choose coordinates on the finite-dimensional image of a linear map.
\writeup Write a tensor as $t=\sum_{i=1}^r\varphi_i\otimes w_i$. Its associated map has image in the span of the finitely many $w_i$, so has finite rank. Choose a basis $u_1,\ldots,u_s$ of that span and rewrite $t=\sum_{j=1}^s\psi_j\otimes u_j$. If its map is zero, then for every $v\in V$,
\[
\sum_j\psi_j(v)u_j=0.
\]
Independence of the $u_j$ gives $\psi_j(v)=0$ for every $j,v$. Hence all $\psi_j$ vanish and $t=0$. This proves injectivity without any dimension restriction on $V$.

If $T\colon V\to W$ has finite-dimensional image with basis $u_1,\ldots,u_s$, let $\lambda_j$ be its coordinate functionals and put $\psi_j=\lambda_jT\in V^*$. Then $T(v)=\sum_j\psi_j(v)u_j$, so $T$ is the image of $\sum_j\psi_j\otimes u_j$. This proves the image assertion.

The identity on $\bigoplus_{j\ge1}Ke_j$ has infinite-dimensional image, so it cannot come from the tensor product. If either the domain or codomain is finite-dimensional, every linear map has finite rank and the map is onto. If either space is zero, both the tensor product and the Hom-space are zero, so it is still an isomorphism.
\end{solution}


\begin{solution}{ex:v10exp-ch03-presentation}
\idea Change the bases of generators and relations until the two relations separate.
\writeup The integer operations give
\[
\begin{pmatrix}4&2\\0&6\end{pmatrix}
\longrightarrow\begin{pmatrix}2&4\\6&0\end{pmatrix}
\longrightarrow\begin{pmatrix}2&4\\0&-12\end{pmatrix}
\longrightarrow\begin{pmatrix}2&0\\0&-12\end{pmatrix}
\longrightarrow\begin{pmatrix}2&0\\0&12\end{pmatrix}.
\]
The steps are a column swap, row subtraction, subtraction of twice the first column from the second, and a row sign change. Each is invertible over $\Z$, so it preserves the isomorphism type of the presented cokernel. Thus $M\cong\Z/2\oplus\Z/12$.

Tensor distributes over this finite direct sum, and the cyclic formula gives $M\otimes\Z/6\cong\Z/2\oplus\Z/6$. A map out of a finite direct sum is a pair of maps. A map from $\Z/r$ to $\Z/6$ is determined by the image of its generator, constrained to be killed by $r$. The two Hom-groups are therefore $\Z/2$ and $\Z/6$. Hence $\operatorname{Hom}(M,\Z/6)\cong\Z/2\oplus\Z/6$, containing twelve homomorphisms. The isomorphism of these two answers is a calculation for this example, not an identification of the tensor and Hom functors.
\end{solution}

\begin{solution}{ex:v10exp-ch03-projective-splitting}
\idea Splitting every surjection onto $P$ is enough, because an arbitrary lifting question can be turned into such a surjection.
\writeup If $P$ is projective and $B\to P$ is onto, lift $\mathrm{id}_P$ through it. The lift is a section, and the splitting lemma applies.

Conversely, suppose all short exact sequences ending in $P$ split. Given $q:B\to C$ onto and $u:P\to C$, the indicated $E$ is a submodule: the equation $u(p)=q(b)$ survives sums and scalar multiples. The projection $\pi:E\to P$ is surjective since each $u(p)$ has a preimage under $q$. Its kernel consists of $(0,b)$ with $b\in\ker q$. Thus
\[
0\to\ker q\xrightarrow{b\mapsto(0,b)}E\xrightarrow{\pi}P\to0
\]
is short exact and has a linear section $s:P\to E$. Since $\pi s=\mathrm{id}_P$, write $s(p)=(p,v(p))$ for a linear map $v:P\to B$. Membership in $E$ gives $q(v(p))=u(p)$. This is the required lift for every $q,u$, proving projectivity. The argument allows $P=0$: its unique map supplies the section and lift.
\end{solution}

\begin{solution}{ex:v10exp-ch03-pid-hypothesis}
\idea Torsion-freeness comes from the ambient domain, while freeness would force a single polynomial generator.
\writeup The ideal is generated by $X,Y$. Since it is contained in the domain $R$, $rf=0$ with $r\ne0$ forces $f=0$, so it is torsion-free. Let $F=K(X,Y)$ be the fraction field. Localising $I$ at nonzero elements gives all of $F$, because $X\in I$ is nonzero and every $z\in F$ is the scalar multiple $(z/X)X$. Hence $F\otimes_RI\cong F$.

If $I$ were free, its basis vectors would remain a basis after this localisation: tensoring the direct sum of copies of $R$ gives the corresponding direct sum of copies of $F$. Its basis would therefore have one member. Thus $I=gR$ for some polynomial $g$. Then $g$ divides both $X$ and $Y$. In the UFD $K[X,Y]$, these are nonassociate irreducibles, so any common divisor is a unit. That would give $I=R$, impossible because evaluation at $(0,0)$ kills $I$ but not $1$. Thus $I$ is not free. The PID hypothesis in Theorem~\ref{thm:pidflat} fails; being a UFD is insufficient.
\end{solution}

\begin{solution}{ex:v10exp-ch03-extensions}
\idea A section must choose a lift of the quotient generator that satisfies its order relation.
\writeup In the relation-lattice description, if $(r,0)=u(n,0)+v(-t,k)$, the second coordinate gives $vk=0$ in $\Z$, hence $v=0$ and $r=un$. Thus $a$ has exactly order $n$. The proposed quotient map is well defined and onto. If an element $ra+sb$ maps to zero, then $s=lk$ and $ra+sb=(r+lt)a$, so its kernel is exactly the cyclic subgroup generated by $a$. This proves short exactness.

Every element mapping to $\bar1$ is uniquely $b+sa$ with $s$ taken modulo $n$. It defines a homomorphism from $\Z/k$ precisely when
\[
k(b+sa)=(t+ks)a=0,
\]
that is, $ks\equiv-t\pmod n$. Put $d=\gcd(k,n)$. Bézout's identity says this congruence is solvable exactly when $d\mid t$. When solvable, writing $k=dk_0$, $n=dn_0$ reduces it to a unique residue class modulo $n_0$, since $k_0$ is invertible there. There are exactly $d$ such classes modulo $n$, hence exactly $d$ sections. A section is a splitting by the splitting lemma.

For $n=k=2,t=0$, the independent relations give $E_0\cong\Z/2\oplus\Z/2$ and there are two sections. For $t=1$, the relation $2b=a$ makes $b$ have order four, so $E_1\cong\Z/4$; the odd right side of $2s\equiv-1\pmod2$ makes splitting impossible. These two middle groups have the same kernel and quotient, but the relation linking the generators changes the lifting problem.
\end{solution}

\section*{Sources and further reading}
\addcontentsline{toc}{section}{Sources and further reading}

The scene of \emph{It's My Turn} (Columbia Pictures, 1980, directed by Claudia Weill) and Benedict Gross's part in it, coaching Clayburgh in her lines, are described in ``Lights, camera and algebraic topology. Imagine'', \emph{The Harvard Crimson}, 23~October~2003; Gross recalls his consulting for the film in R.~Kirby's interview on Celebratio Mathematica (\texttt{celebratio.org/Gross\_BH}).
H.~Whitney, ``Tensor products of Abelian groups'', \emph{Duke Math.\ J.} 4 (1938), 495--528; R.~Baer, ``Abelian groups that are direct summands of every containing Abelian group'', \emph{Bull.\ Amer.\ Math.\ Soc.} 46 (1940),
800--806; V.~Strassen, ``Gaussian elimination is not optimal'', \emph{Numer.\ Math.} 13 (1969), 354--356; S.~Winograd, ``On multiplication of $2\times2$ matrices'', \emph{Linear Algebra Appl.} 4 (1971), 381--388.

For more: M.~F.~Atiyah and I.~G.~Macdonald, \emph{Introduction to Commutative Algebra} (Addison-Wesley, 1969), Chapter~2, which covers this chapter in twenty pages; S.~Lang, \emph{Algebra} (Springer, 3rd ed.\ 2002),
Chapters~III and XVI; and J.~J.~Rotman, \emph{An Introduction to Homological Algebra} (Springer, 2nd ed.\ 2009), Chapters~2 and~3, with everything about projective, injective and flat modules.
